% Compile with XeLaTeX, LuaLaTeX, or Tectonic for multilingual text.
\documentclass[12pt,a4paper]{report}
\usepackage{iftex}
\RequireTUTeX
\usepackage[a4paper,left=30mm,right=25mm,top=25mm,bottom=25mm]{geometry}
\usepackage{fontspec}
\setmainfont{texgyretermes-regular.otf}[
  BoldFont=texgyretermes-bold.otf,
  ItalicFont=texgyretermes-italic.otf,
  BoldItalicFont=texgyretermes-bolditalic.otf]
\setsansfont{texgyreheros-regular.otf}[
  BoldFont=texgyreheros-bold.otf,
  ItalicFont=texgyreheros-italic.otf]
\setmonofont{texgyrecursor-regular.otf}[
  BoldFont=texgyrecursor-bold.otf,
  ItalicFont=texgyrecursor-italic.otf]
\usepackage{amsthm,amssymb,mathrsfs,setspace}%amsmath, latexsym,footmisc

\usepackage[nottoc,notlof,notlot]{tocbibind}
\renewcommand{\chaptermark}[1]{\markboth{#1}{}}
\renewcommand{\sectionmark}[1]{\markright{\thesection\ #1}}

% ── Additional packages for thesis content ───────────────────────
\usepackage{amsmath}
\usepackage{graphicx}
\usepackage{float}
\usepackage{caption}
\captionsetup{font=small,labelfont=bf,textfont=it,
              justification=centering,skip=6pt}
\usepackage{booktabs,array,longtable,tabularx,multirow}
\usepackage[protrusion=true,expansion=false]{microtype}
\usepackage{seqsplit}
\renewcommand{\arraystretch}{1.3}
\usepackage{xcolor}
\usepackage{listings}
\lstset{basicstyle=\small\ttfamily,breaklines=true,
        frame=single,backgroundcolor=\color{gray!8},
        showstringspaces=false}
\usepackage{tikz}
\usetikzlibrary{positioning,fit,arrows.meta}
\usepackage{enumitem}
\usepackage[hidelinks,
  pdftitle={Multilingual Conversational AI for Labour Force Surveys},
  pdfauthor={Sivarama Krishnan Subbiah}]{hyperref}
\usepackage{xurl}
\urlstyle{same}
\usepackage[numbers,sort&compress]{natbib}
% The thesis prose is English; application language support is documented separately.
\usepackage{polyglossia}
\setdefaultlanguage{english}
\setlength{\emergencystretch}{3em}

% ── Theorem environments ─────────────────────────────────────────
\theoremstyle{plain}
\newtheorem{theorem}{Theorem}[section]
\newtheorem{lemma}[theorem]{Lemma}
\newtheorem{corollary}[theorem]{Corollary}
\newtheorem{proposition}[theorem]{Proposition}
\theoremstyle{definition}
\newtheorem{definition}[theorem]{Definition}
\newtheorem{example}[theorem]{Example}
\newtheorem{notation}[theorem]{Notation}
\theoremstyle{remark}
\newtheorem{remark}[theorem]{Remark}

\onehalfspacing
\setlength{\headheight}{15pt}

\begin{document}
\hypersetup{pageanchor=false}

% --------------- Title page -----------------------
\begin{titlepage}
\singlespacing

\begin{center}

\vspace*{-0.5cm}

\textbf{\Large MULTILINGUAL CONVERSATIONAL AI FOR LABOUR FORCE SURVEYS:
A MULTI-AGENT RAG SYSTEM WITH CREWAI FRAMEWORK}\\[10pt]

\vspace*{1.7cm}

                         A Project Report Submitted \\
                     in Partial Fulfilment of the Requirements  \\
                     for the Degree of  \\
                          \vspace{3mm}
                   {\Large \bf Master of Technology } \\
                   in \\
                   {\large \bf Artificial Intelligence and Data Science } \\
                      \vspace{5mm}
                   {\em  by} \\ \vspace{2mm}
             {\large \bf Sivarama Krishnan Subbiah} \\
{\large \bf (Roll No. 2024MCS110007)}\\[0.2in]

\vfill

\includegraphics[width=5cm,keepaspectratio]{logo2.jpg}
\vspace*{0.4cm}

{\em\large to }\\
{\bf\large INDIAN INSTITUTE OF INFORMATION TECHNOLOGY}\\
{\bf\large KOTTAYAM-686635, INDIA}\\
{\it\large 2026}
\end{center}

\end{titlepage}
\hypersetup{pageanchor=true}

\clearpage
% --------------- Declaration page -----------------------
\pagenumbering{roman} \setcounter{page}{2}
\begin{center}
{\large{\bf{DECLARATION}}}
\end{center}

\noindent I, \textbf{Sivarama Krishnan Subbiah} (\textbf{Roll No: 2024MCS110007}),
hereby declare that this report entitled
\textbf{``Multilingual Conversational AI for Labour Force Surveys:
A Multi-Agent RAG System with CrewAI Framework''} submitted to Indian
Institute of Information Technology Kottayam towards partial
requirement of {\bf Master of Technology} in \textbf{Artificial Intelligence
and Data Science}
is an original work carried out by me under the supervision of
\textbf{Dr.\ Goutam Mali} and has not formed the basis for the award
of any degree or diploma, in this or any other institution or
university. External information, methods, and results used in this
report have been acknowledged and cited. The reported experiments are
distinguished from planned studies and unverified claims.

\vspace{4cm}

\noindent Kottayam-686635 \hfill \textbf{Sivarama Krishnan Subbiah}

\noindent 2026

\clearpage

% --------------- Certificate page -----------------------
\pagenumbering{roman} \setcounter{page}{3}
\begin{center}
{\large{\bf{CERTIFICATE}}}
\end{center}

\noindent This is to certify that the work contained in this
project report entitled \textbf{``Multilingual Conversational AI for
Labour Force Surveys: A Multi-Agent RAG System with CrewAI Framework''}
submitted by \textbf{Sivarama Krishnan Subbiah}
(\textbf{Roll No: 2024MCS110007}) to Indian Institute of Information
Technology Kottayam towards partial requirement of {\bf Master of Technology}
in \textbf{Artificial Intelligence and Data Science}
has been carried out by him under my supervision and that it has not been
submitted elsewhere for the award of any degree.

\vspace{4cm}

\noindent Kottayam-686635  \hfill (Dr.\ Goutam Mali)

\noindent 2026 \hfill Project Supervisor

\clearpage

% --------------- Abstract page -----------------------
\begin{center}
{\Large{\bf{ABSTRACT}}}
\end{center}

Labour Force Surveys require consistent measurement of employment and
standardised coding of occupations, industries, and education. Multilingual
respondent descriptions and ambiguous job titles make automated coding
difficult. This thesis develops and evaluates a conversational prototype
targeting English, Arabic, Urdu, Hindi, and Tagalog, with dictionary-based
Gulf Arabic normalisation and a modular CrewAI-assisted backend.

The system combines a deterministic interview state machine with
retrieval-based ISCO-08 classification, ISIC Rev.~4 industry classification,
and ISCED~2011 attainment and ISCED-F~2013 field mapping. A Semantic Relation
Engine checks occupation--industry and occupation--attainment compatibility
using disclosed, project-authored rules. Its coherence score is a heuristic,
and HIGH violations trigger human-review queue creation. Rule and endpoint
tests establish agreement with implementation expectations; independent
expert validation and probability calibration remain outstanding.

On a project-generated, occupation-group-separated WISCO benchmark of
18,747 language-specific reference titles, flat retrieval achieved
21.19\% exact unit-group accuracy compared with 10.35\% for hierarchical
retrieval. Combining official-definition enrichment with a larger intended
multilingual encoder produced the highest observed accuracy, 40.95\%;
that run's preserved records do not establish which encoder actually
executed, so the gain is not attributed to encoder size. The configuration
the prototype actually serves is a fragment-to-parent retriever reaching
38.83\% on the same records, 6.28 points above its own same-encoder
baseline and bound to a recorded serving identity.
These later comparisons reused the same benchmark during development and
are therefore interpreted as comparative evidence requiring confirmation
on an untouched test set. Tested LLM reranking configurations did not show
consistent additional gains. Industry and education-field experiments use
synthetic data and are reported separately from occupation-reference data
and human survey evidence.

The thesis contributes a comparison of retrieval designs with preserved predictions,
an explicit cross-standard review mechanism, and an analysis of the
prototype's empirical and operational limitations. A synthetic
30-scenario interview run demonstrates execution of the implemented
workflow, but does not measure respondent usability or burden reduction.
The planned human pilot has not been conducted and requires ethics
approval before recruitment. Routing, persistence, and access-control
defects identified in review have been repaired and regression-tested,
with re-verification in the deployment environment still outstanding. The results
support further research on multilingual classification and assisted
survey coding rather than immediate replacement of a validated field
survey instrument.

\vspace{6pt}
\noindent\textbf{Keywords:} Labour Force Survey; conversational AI;
multi-agent systems; CrewAI; retrieval-augmented generation; ISCO-08;
ISIC Rev.4; ISCED-F 2013; Semantic Relation Engine; multilingual NLP;
code-switching; human-in-the-loop; official statistics.

\clearpage


\tableofcontents
\clearpage
\listoffigures
\listoftables


\newpage

\pagenumbering{arabic}
\setcounter{page}{1}

% ====== Main chapters ===========================================

\input chapter1.tex

\input chapter2.tex

\input chapter3.tex

\input chapter4.tex

\input chapter5.tex

\input chapter6.tex

\input chapter7.tex

\input appendices.tex

\bibliographystyle{plainnat}
\bibliography{bib}

\end{document}
